\SourceLeaf{040}
\fancyhead[L]{\small Another New Proof}
\begin{center}
\textbf{ANOTHER NEW PROOF OF THE THEOREM}\\[12pt]
\textbf{THAT EVERY POLYNOMIAL}\\[12pt]
IN ONE VARIABLE\\[12pt]
\textbf{\parbox{.95\linewidth}{\centering CAN BE RESOLVED INTO REAL FACTORS\\OF THE FIRST OR SECOND DEGREE.}}\\[12pt]
\rule{3cm}{0.3pt}
\end{center}

\SecondArticle{1} Although the proof of the theorem on resolving polynomials into factors that I gave in a paper published sixteen years ago seems to leave nothing to be desired either in rigour or in simplicity, I hope that mathematicians will not find it unwelcome if I return to this most important question and attempt to establish another proof, no less rigorous, from entirely different principles. That earlier proof depends, at least in part, on geometric considerations; the one I now undertake to set out, by contrast, will rest on purely analytic principles. In the work cited, I reviewed the more notable analytic methods by which other mathematicians had tried to prove our theorem up to that time, and gave a detailed account of the defects they suffer from. The most serious and truly fundamental of these is common to all those attempts, as well as to the more recent ones that have come to my attention; yet even then I declared that it seemed by no means unavoidable in an analytic proof. Let the experts now judge whether these further efforts have fully redeemed the pledge I gave at that time.

\SecondArticle{2} Some preliminary investigations will precede the main discussion, both so that nothing may seem to be lacking and because their treatment may perhaps shed some fresh light even on matters that others have already\SourceLeaf{041}
 touched upon. We shall first consider the greatest common divisor of two polynomials in one indeterminate. It should be stated at the outset that we are always speaking here only of polynomials: if a product is formed from two of them, each is called a \textit{divisor} of that product. The \textit{degree} of a divisor is determined by the exponent of the highest power of the indeterminate it contains, without any regard at all to the numerical coefficients. We can deal with the remaining matters concerning common divisors of polynomials more briefly, since they are entirely analogous to the corresponding properties of common divisors of numbers.

Given two polynomials $Y,Y'$ in the indeterminate $x$, of which the former has a higher degree than the latter, or at least not a lower one, we shall form the following equations:
\[
\begin{gathered} Y=qY'+Y''\\ Y'=q'Y''+Y'''\\ Y''=q''Y'''+Y''''\\ \text{etc., up to}\\ Y^{(\mu-1)}=q^{(\mu-1)}Y^{(\mu)}
\end{gathered}
\]
The rule is to divide $Y$ first by $Y'$ in the usual way; then $Y'$ by the remainder $Y''$ of the first division, which has lower degree than $Y'$; then again the first remainder by the second, $Y'''$; and so on until a division without remainder is reached. That this must eventually happen follows from the continually decreasing degrees of the polynomials $Y',Y'',Y'''$ etc. There is hardly any need to point out that these functions, as well as the quotients $q,q',q''$ etc., are \textit{polynomials} in $x$. With this established, it is clear that:

I. Going backwards from the last of these equations to the first, the polynomial $Y^{(\mu)}$ is seen to divide each of the preceding ones, and is therefore certainly a common divisor of the given polynomials $Y,Y'$.

II. Going forwards from the first equation to the last, we see that every common divisor of the polynomials $Y,Y'$ also divides each of the following ones, and hence the last, $Y^{(\mu)}$. Thus the polynomials $Y,Y'$ cannot have any common divisor of higher degree than $Y^{(\mu)}$. Every common divisor of the same degree as $Y^{(\mu)}$ will stand to it in the ratio of a number to a number; hence it must itself be taken as the greatest common divisor.
\par
\SourceLeaf{042}
III. If $Y^{(\mu)}$ has degree 0, that is, is a number, no function of the indeterminate $x$ in the proper sense can divide $Y,Y'$. In this case, therefore, these polynomials are said to have no common divisor.

\Needspace{18\baselineskip}IV. Take the equation immediately before the last one; then eliminate $Y^{(\mu-1)}$ from it with the aid of the equation immediately before that; then again eliminate $Y^{(\mu-2)}$ with the aid of the preceding equation, and so on. In this way we obtain
\[
\begin{aligned}
Y^{(\mu)}&=+kY^{(\mu-2)}-k'Y^{(\mu-1)}\\
&=-k'Y^{(\mu-3)}+k''Y^{(\mu-2)}\\
&=+k''Y^{(\mu-4)}-k'''Y^{(\mu-3)}\\
&=-k'''Y^{(\mu-5)}+k''''Y^{(\mu-4)}\\
&\quad\text{etc.}
\end{aligned}
\]
if we suppose the polynomials $k,k',k''$ etc. to be formed by the following rule:
\[
\begin{aligned}
k&=1\\
k'&=q^{(\mu-2)}\\
k''&=q^{(\mu-3)}k'+k\\
k'''&=q^{(\mu-4)}k''+k'\\
k''''&=q^{(\mu-5)}k'''+k''\\
&\quad\text{etc.}
\end{aligned}
\]
We shall therefore have
\[
\pm k^{(\mu-2)}Y\mp k^{(\mu-1)}Y'=Y^{(\mu)}
\]
with the upper signs applying when $\mu$ is even and the lower when it is odd. Thus, when $Y$ and $Y'$ have no common divisor, we can in this way find two polynomials $Z,Z'$ in the indeterminate $x$ such that
\[
ZY+Z'Y'=1
\]
V. The converse of this proposition clearly holds as well: if the equation
\[
ZY+Z'Y'=1
\]
can be satisfied with $Z,Z'$ polynomials in the indeterminate $x$, then $Y$ and $Y'$ certainly cannot have a common divisor.

\SourceLeaf{043}
\Needspace{16\baselineskip}\SecondArticle{3}
The second preliminary investigation will concern the transformation of symmetric functions. Let $a,b,c$ etc. be indeterminate quantities, $m$ in number. Denote their sum by $\lambda'$, the sum of their products in pairs by $\lambda''$, the sum of their products in threes by $\lambda'''$, etc., so that expanding the product
\[
(x-a)(x-b)(x-c)\, .\, .\, .\, .
\]
gives
\[
x^m-\lambda'x^{m-1}+\lambda''x^{m-2}-\lambda'''x^{m-3}+\text{etc.}
\]
Thus $\lambda',\lambda'',\lambda'''$ etc. are symmetric functions of the indeterminates $a,b,c$ etc.: that is, functions in which these indeterminates occur in the same way, or more explicitly, functions unchanged by any permutation of these indeterminates among themselves. More generally, every polynomial in $\lambda',\lambda'',\lambda'''$ etc. (whether it contains only these indeterminates or also others independent of $a,b,c$ etc.) will clearly be a symmetric polynomial in the indeterminates $a,b,c$ etc.

\SecondArticle{4}
The converse theorem is somewhat less evident. Let $\rho$ be a symmetric function of the indeterminates $a,b,c$ etc.; it will thus consist of a certain number of terms of the form
\[
Ma^\alpha b^\beta c^\gamma\, .\, .\, .\, .
\]
where $\alpha,\beta,\gamma$ etc. denote nonnegative integers, and $M$ is a coefficient either fixed or at least independent of $a,b,c$ etc. (if the function $\rho$ happens to contain other indeterminates besides $a,b,c$ etc.). First we shall establish a definite ordering of these individual terms. For this purpose, we first arrange the indeterminates $a,b,c$ etc. in a definite order, which is itself entirely arbitrary: for example, with $a$ in the first place, $b$ in the second, $c$ in the third, etc. Then, of the two terms
\[
Ma^\alpha b^\beta c^\gamma\, .\, .\, .\, .\quad\text{and}\quad Ma^{\alpha'}b^{\beta'}c^{\gamma'}\, .\, .\, .\, .
\]
we shall place the former higher in the ordering than the latter if
\[
\begin{gathered}\text{or}\quad\alpha>\alpha',\quad\text{or}\quad\alpha=\alpha'\text{ and }\beta>\beta',\\\text{or}\quad\alpha=\alpha',\ \beta=\beta'\text{ and }\gamma>\gamma',\quad\text{or etc.}\end{gathered}
\]

\SourceLeaf{044}
that is, if the first nonvanishing one among the differences $\alpha-\alpha',\beta-\beta',\gamma-\gamma'$ etc. is positive. Terms in the same place in this ordering therefore differ only in their coefficient $M$, and can be combined into one term. We shall accordingly suppose the individual terms of the function $\rho$ to occupy different places in the ordering.

Now observe that, if $Ma^\alpha b^\beta c^\gamma\, .\, .\, .\, .\, .$ is the highest of all the terms of the function $\rho$ in this ordering, $\alpha$ must be greater than, or at least not less than, $\beta$. For if $\beta>\alpha$, the term $Ma^\beta b^\alpha c^\gamma\, .\, .\, .\, .\, .$, which the function $\rho$ must also contain since it is symmetric, would stand higher than $Ma^\alpha b^\beta c^\gamma\, .\, .\, .\, .\, .$, contrary to the hypothesis. In the same way, $\beta$ will be greater than, or at least not less than, $\gamma$; further, $\gamma$ will not be less than the next exponent $\delta$, etc. Hence each of the differences $\alpha-\beta$, $\beta-\gamma,\gamma-\delta$ etc. is a nonnegative integer.

\Needspace{15\baselineskip}Secondly, consider a product formed from any number of polynomials in the indeterminates $a,b,c$ etc. Its highest term must be the product of the highest terms of those factors. It is equally clear that the highest terms of the functions $\lambda',\lambda'',\lambda'''$ etc. are respectively $a,ab,abc$ etc. It follows that the highest term arising from the product
\[
p=M{\lambda'}^{\alpha-\beta}{\lambda''}^{\beta-\gamma}{\lambda'''}^{\gamma-\delta}\, .\, .\, .\, .\, .
\]
is $Ma^\alpha b^\beta c^\gamma\, .\, .\, .\, .\, .$. Thus, on setting $\rho-p=\rho'$, the highest term of the function $\rho'$ will certainly stand lower in the ordering than the highest term of the function $\rho$. But $p$, and therefore also $\rho'$, are clearly symmetric polynomials in $a,b,c$ etc. Treating $\rho'$ in the same way as $\rho$ before it, we can therefore split it into $p'+\rho''$, where $p'$ is a product of powers of $\lambda',\lambda'',\lambda'''$ etc. multiplied by a coefficient either fixed or at least independent of $a,b,c$ etc., while $\rho''$ is a symmetric polynomial in $a,b,c$ etc. whose highest term stands lower in the ordering than the highest term of $\rho'$. Continuing in the same way, $\rho$ will clearly at last be reduced to the form $p+p'+p''+p'''$ etc., that is, transformed into a polynomial in $\lambda',\lambda'',\lambda'''$ etc.

\SecondArticle{5}
We can also state the theorem proved in the preceding article as follows. Given any symmetric polynomial $\rho$ in the indeterminates $a,b,c$ etc., we can specify a polynomial in as many other indeterminates $l',l'',l'''$ etc. which becomes
\SourceLeaf{045}
$\rho$ under the substitutions $l'=\lambda',l''=\lambda'',l'''=\lambda'''$ etc. It is also easy to show that \textit{this can be done in only one way.} For suppose that two different functions of the indeterminates $l',l'',l'''$ etc., namely $r$ and $r'$, yield the same function of $a,b,c$ etc. after the substitutions $l'=\lambda',l''=\lambda'',l'''=\lambda'''$ etc. Then $r-r'$ will be a function of $l',l'',l'''$ etc. which does not vanish in itself, but vanishes identically after those substitutions. We shall readily see that this is absurd if we consider that $r-r'$ must consist of a certain number of parts of the form
\[
M{l'}^\alpha{l''}^\beta{l'''}^\gamma\, .\, .\, .\, .\, .
\]
whose coefficients $M$ do not vanish and which all differ from one another in their exponents. \Needspace{8\baselineskip}The highest terms arising from these individual parts are therefore given by
\[
Ma^{\alpha+\beta+\gamma+\text{etc.}}b^{\beta+\gamma+\text{etc.}}c^{\gamma+\text{etc.}}\, .\, .\, .\, .\, .
\]
and must consequently occupy different places in the ordering, so that the highest term of all cannot possibly cancel.

The actual calculation for transformations of this kind could, moreover, be greatly shortened by various shortcuts. We do not dwell on them here, since the mere possibility of the transformation already suffices for our purpose.
